\documentclass[11pt,twoside]{amsart}
\usepackage{amssymb, amsmath, enumerate, palatino, hyperref}
\usepackage[normalem]{ulem}
\usepackage{fullpage}
\usepackage[T1]{fontenc}
\renewcommand{\labelitemi}{\guillemotright}
\usepackage{mathrsfs}

\usepackage{tikz}
\tikzstyle{ball} = [circle,shading=ball, ball color=black,
    minimum size=1mm,inner sep=1.3pt]

\theoremstyle{plain}
\newtheorem{prop}{Proposition}%[section]
\newtheorem{lemma}[prop]{Lemma}
\newtheorem{thm}[prop]{Theorem}
\newtheorem{obs}[prop]{Observation}
\newtheorem{app}[prop]{Application}
\newtheorem*{MainThm}{Main Theorem}
\newtheorem{cor}[prop]{Corollary}
\newtheorem{conj}[prop]{Conjecture}
\theoremstyle{remark}
\newtheorem{rmk}[prop]{Remark}
\newtheorem{prob}{Problem}
\newtheorem{bonus}[prop]{Bonus Problem}
\newtheorem{exc}{Exercise}
\theoremstyle{definition}
\newtheorem{ex}[prop]{Example}
\theoremstyle{definition}
\newtheorem{defn}[prop]{Definition}

\newcommand{\RR}{\mathbb{R}}
\newcommand{\ZZ}{\mathbb{Z}}
\newcommand{\CC}{\mathbb{C}}
\newcommand{\NN}{\mathbb{N}}
\newcommand{\QQ}{\mathbb{Q}}
\newcommand{\PP}{\mathbb{P}}

\newcommand{\cC}{\mathscr{C}}
\newcommand{\Ob}{\operatorname{Ob}}
\newcommand{\Mor}{\operatorname{Mor}}

\newcommand{\Set}{\operatorname{Set}}
\newcommand{\FinSet}{\operatorname{FinSet}}
\newcommand{\Vect}{\operatorname{Vect}}
\newcommand{\FinVect}{\operatorname{FinVect}}
\newcommand{\Mat}{\operatorname{Mat}}
\newcommand{\Gp}{\operatorname{Gp}}
\newcommand{\FinGp}{\operatorname{FinGp}}
\newcommand{\AbGp}{\operatorname{AbGp}}
\newcommand{\Ring}{\operatorname{Ring}}
\newcommand{\CommRing}{\operatorname{CommRing}}
\newcommand{\Field}{\operatorname{Field}}
\newcommand{\Top}{\operatorname{Top}}
\newcommand{\Aut}{\operatorname{Aut}}
\newcommand{\id}{\operatorname{id}}
\DeclareRobustCommand{\stir}{\genfrac\{\}{0pt}{}}

% For solutions.
\newenvironment{sol}{
  \medskip

  \noindent\textsc{Solution:}
}
{
  \medskip

}

\newcommand{\defeq}{:=}

\title{Math 113: Discrete Structures\\ Homework 33}
%\author{} %Remove the leading % and put your name here if using this .tex file as a template for your homework.

\begin{document}
\maketitle




\begin{prob}  Use the Euclidean algorithm to compute the following (showing your work):
\[
  \text{(a)}\quad
  \gcd(28,63)\qquad\qquad\text{(b)}\quad\gcd(234,286)\qquad\qquad\text{(c)}\quad\gcd(51,894).
\]
\end{prob}
\medskip

\begin{prob}
Use the Euclidean algorithm to compute the $\gcd$ of $204$ and $168$ and then
use back-substitution to find
integers $m$ and~$n$ such that
\[
  \gcd(204,168) = 204m+168n.
\]
Show your work.  Remember to use back-substitution and not the extended
Euclidean algorithm.
\end{prob}
\medskip

\begin{prob} 
  \ 
  \begin{enumerate}[(a)]
    \item Show that if~$n$ is positive integer of the form~$4k+3$ for some
      integer~$k$, then~$n$ is not a perfect square. (Hint: Suppose~$n=m^2$.  We
      can then write~$m=4q+r$ for some~$r\in\left\{ 0,1,2,3 \right\}$.  Consider
    the remainders of the quantities~$(4q)^2$, $(4q+1)^2$,
    $(4q+2)^2$, and $(4q+3)^2$ upon division by~$4$.)
  \item Show that no integer in the sequence
    \[
      11,\quad 111,\quad 1111,\quad 11111,\quad \dots
    \]
    is a perfect square. 
  \end{enumerate}
\end{prob}
\end{document}
